\documentclass[11pt,a4paper]{article}
% Shared preamble for RuPhO English problem statements (match T1 2026 style).
% Use: \input{latex/rupho_en_preamble.tex} after \documentclass.
\usepackage[margin=1in]{geometry}
\usepackage{amsmath,amssymb}
\usepackage{graphicx}
\usepackage{float}
\usepackage{enumitem}
\usepackage{microtype}
\usepackage{caption}
\usepackage{tabularx}
\usepackage{hyperref}

\captionsetup{font=small,labelfont=bf,skip=6pt}

\setlist[itemize]{leftmargin=1.2cm}
\setlist[enumerate]{leftmargin=1.2cm}

% One outer box per subproblem: [id | statement | points] (no inner rules)
\newcommand{\problembox}[3]{%
  \par\addvspace{0.42em}%
  \noindent
  \setlength{\fboxsep}{7.5pt}%
  \setlength{\fboxrule}{0.95pt}%
  \fbox{%
    \begin{minipage}{\dimexpr\linewidth-2\fboxsep-2\fboxrule\relax}
    \begin{tabularx}{\linewidth}{@{}>{\raggedright\arraybackslash\bfseries}p{2.1em} @{\hspace{0.9em}} >{\raggedright\arraybackslash}X @{\hspace{0.9em}} >{\raggedleft\arraybackslash\bfseries}p{2.4em}@{}}
    #1 & #2 & #3
    \end{tabularx}
    \end{minipage}%
  }%
  \par\addvspace{0.32em}%
}


\title{\textbf{Superconducting films}}
\author{\normalsize X21 (2021) --- English translation}
\date{}
\begin{document}
\maketitle

In thin films $\rm In-O$ at low temperatures (of the order of $1~\text{K}$), a superconductor-insulator phase transition controlled by an external magnetic field is observed (the transition occurs when the external magnetic field $B$ changes).

\section*{Part A. Critical parameters (6.8 points)}

The table “Experimental data” presents measurements of the dependences $R(T)$, where $R$ is the resistance of the sample, $T$ is the temperature, at different external magnetic fields $B$. Resistance measurement error $\varepsilon(R)=0.5\%$. The graph (see Fig.) shows approximate types of dependencies $R(T)$ for some values $B$ indicated in the table. Let us assume that the dependence of resistance $R(T,B)$ has the following form: $$R(T,B) = R_C - \gamma T + A \left( B - B_C \right) T^\{\alpha\}, $$ где $A, R_C, B_C$ — некоторые параметры, связанные со свойствами образца. Если на кривой $R(T)$ есть точка перегиба (то есть $\dfrac{\partial R}{\partial T}=0$), то обозначим ее координаты, как $\left(T_\text{п}, R_\text{n}\right)$. Примечание: Для функции двух переменных аналогом производной выступает частная производная согласно определениям:   Фактически, при вычислении частой производной по $x$ мы полагаем $y$ constant and vice versa, for example: \\[ f(x,y) = 3x^2 + 2xy,\quad \frac\{\partial\}\{\partial x\}f(x,y) = 6x + 2y,\quad \frac\{\partial\}\{\partial y\}f(x,y) = 2x \\]

\begin{figure}[H]
    \centering
    \includegraphics[width=0.72\textwidth]{PE_EN_assets/fig1.jpg}
    \caption{Figure 1.}
\end{figure}

\problembox{A1}{Derive the analytical relationship $R_\text{п}$ ($T_\text{п}$).}{1.00}

\problembox{A2}{Define $B_C$.}{0.40}

\problembox{A3}{Define $R_C$ in two ways.}{2.60}

\problembox{A4}{Define the parameters $\gamma$ and $\alpha$.}{2.80}

\section*{Part B. Probability of describing an experiment (3.2 points)}

Estimating the error of a physical quantity $X$ allows one to calculate the probability with which its value lies in the interval $[a,b]$. To assess the probability $p$ that the result of the experiment is described by a mathematical expression, the criterion $\chi^2$ (chi-square) is used, calculated according to the following formula: \\[ \chi^2 = \sum\limits_\{i=1\}^\{N\} \left( \frac\{R_\text\{теор\} - R_\text\{эксп\}\}\{\Delta R\} \right)^2 \\] For a number of points $N=150$, a correspondence table $p$ and $\chi^2$:

(Probability table)

\problembox{B1}{Count $\chi^2$ (definition in note) of $R(T)$ dependencies for each $B$.}{2.20}

\problembox{B2}{Using the table from the note, estimate the probability according to which the formula $f={R(T,B)}$ describes the dependence $R(T,B)$ for any (specify which) $150$ points.}{1.00}

\vspace{1em}
\noindent\textit{Source:} \url{https://pho.rs/p/102}

\end{document}